
\section{Appendix A. Mathematical Derivations}

\subsection{A.1 Prefix Sum Correctness}

\textbf{Claim.} The cumulative sum formulation computes exact causal linear 
attention.

\textbf{Proof.} For position $i$, causal linear attention is:

\[
O_i = \frac{\sum_{j \le i} \phi(Q_i)^\top \phi(K_j) V_j}{\sum_{j \le i} \phi(Q_i)^\top \phi(K_j)}
\]

Factoring out $\phi(Q_i)$:

\[
O_i = \frac{\phi(Q_i)^\top \sum_{j \le i} \phi(K_j) V_j^\top}{\phi(Q_i)^\top \sum_{j \le i} \phi(K_j)}
\]

The cumulative sums $\text{cum}^{(i)} = \sum_{j \le i} \phi(K_j) V_j^\top$ 
and $\text{den}^{(i)} = \sum_{j \le i} \phi(K_j)$ are computed via \texttt{cumsum}. 
$\square$

\subsection{A.2 Non-Negativity of ELU+1}

\textbf{Claim.} $\phi(x) = \mathrm{elu}(x) + 1 > 0$ for all $x \in \mathbb{R}$.

\textbf{Proof.} 

\begin{itemize}
\item For $x > 0$: $\mathrm{elu}(x) = x > 0$, so $\phi(x) = x + 1 > 1 > 0$.
\item For $x \le 0$: $\mathrm{elu}(x) = e^x - 1 \in (-1, 0]$, so
  $\phi(x) = e^x \in (0, 1]$.

\end{itemize}
In both cases, $\phi(x) > 0$. $\square$

\subsection{A.3 Gate Gradient}

The gradient of the loss w.r.t. the gate logit $z = w^\top x + b$:

\[
\frac{\partial \mathcal{L}}{\partial z} = \frac{\partial \mathcal{L}}{\partial g} \cdot g(1 - g)
\]

where:

\[
\frac{\partial \mathcal{L}}{\partial g} = \nabla_O \mathcal{L} \cdot (O^\text{dense} - O^\text{lin})
\]

\subsection{A.4 Fusion Interpolation}

\textbf{Claim.} The fusion $O = g O^\text{dense} + (1 - g) O^\text{lin}$ is a 
convex combination for $g \in [0, 1]$.

\textbf{Proof.} Trivial from $g + (1 - g) = 1$ and $0 \le g \le 1$. $\square$

\subsection{A.5 Cumsum Memory}

For a batch of $B$ sequences of length $L$ with $H$ heads and head dim $d_h$:

\[
\text{Memory} = B H L d_h^2 \cdot \text{bytes}
\]

For $B = 8, H = 12, L = 512, d_h = 64$: $8 \times 12 \times 512 \times 64^2 \times 2 = 4$ GB. 

\textbf{This is prohibitive.} We need chunked cumsum for practical long-context 
inference.

\subsection{A.6 Approximation Error}

\textbf{Claim.} The fusion error is bounded by:

\[
\|O - O^\text{dense}\| \leq (1 - g) \cdot \|O^\text{lin} - O^\text{dense}\|
\]

\textbf{Proof.} $O - O^\text{dense} = (1 - g)(O^\text{lin} - O^\text{dense})$. 
Taking norms gives (A.6). $\square$
\begin{lstlisting}

---
\end{lstlisting}